\section{Titularidade e Propriedade Intelectual}
\label{sec:propriedade-intelectual}

A ontologia do gêmeo digital cria um complexo mosaico jurídico em torno da sua titularidade. Diferentemente de uma radiografia tradicional, que é um arquivo estático de posse do paciente e guarda da clínica, o Gêmeo Digital é um artefato dinâmico construído colaborativamente.

Quem é o dono do modelo?
\begin{itemize}
    \item O \textbf{paciente} detém a titularidade sobre os dados brutos biométricos subjacentes, amparado pela LGPD, possuindo o direito inalienável de portabilidade.
    \item O \textbf{cirurgião-dentista} detém direitos autorais sobre a inteligência clínica (\textit{know-how}) aplicada na correção das malhas, demarcação de margens e configuração vetorial do tratamento.
    \item A \textbf{empresa de software} detém a propriedade intelectual sobre os algoritmos de segmentação, as bibliotecas de materiais in silico e a infraestrutura de processamento computacional \cite{mdpi2025digital}.
\end{itemize}

O vácuo doutrinário sobre este mosaico gera fricções colossais. Quando um paciente decide mudar de dentista, as plataformas proprietárias frequentemente bloqueiam a exportação do arquivo planejado (o \textit{setup} dinâmico), permitindo apenas a exportação de malhas burras (arquivos STL sem metadados biomecânicos). Esta prática (\textit{data hoarding}) atenta contra o direito do consumidor e a continuidade do cuidado. A padronização de contratos de prestação de serviços deve evoluir para delimitar claramente os direitos de co-propriedade sobre este novo ativo.
